\label{cap:avaliacao}

Este capítulo discute os resultados apresentados no Capítulo~5 em relação às questões de pesquisa definidas na Introdução e aos critérios de avaliação descritos na Seção~\ref{sec:delineamento-avaliacao}.

\section{RQ1 -- Aplicabilidade e Limitações}
\label{sec:avaliacao-rq1}

A RQ1 examina a aplicabilidade da abordagem proposta a partir da reconstrução retrospectiva parcial da migração da BSN. A avaliação considera a viabilidade das atividades (C1), a utilidade dos artefatos e resultados para as etapas posteriores (C2) e as limitações identificadas (C3), conforme os critérios estabelecidos na Seção~\ref{sec:delineamento-avaliacao}.

\paragraph{C1 -- Viabilidade das atividades.}

Os resultados do Capítulo~5 mostram que foi possível recuperar e reconstruir evidências correspondentes às atividades da abordagem. A Compreensão de Programas foi apoiada pelas visões estáticas e dinâmicas produzidas em trabalhos anteriores sobre a BSN, enquanto a Especificação Comportamental utilizou propriedades documentadas e cenários BDD disponíveis no legado. O Estabelecimento da \textit{Ground Truth} Comportamental foi examinado a partir dos testes disponíveis e das evidências de sua execução no ROS~1, considerando seus critérios de aceitação e os resultados observados.

Na fase de migração, a Identificação das Realizações dos Comportamentos foi sistematizada retrospectivamente por meio da inspeção do código-fonte, dos grafos de chamadas e das interfaces de comunicação. Essas informações permitiram relacionar cenários BDD com entidades e responsabilidades da implementação legada. A partir desse mapeamento, foram caracterizadas as transformações T1--T3, apresentadas na Seção~\ref{sec:resultados-planejamento-transformacoes}. A implementação existente no ROS~2 forneceu as evidências das mudanças realizadas, enquanto a Validação Comportamental foi examinada pela comparação dos testes, oráculos e resultados disponíveis nas duas versões.

Entretanto, a viabilidade observada não corresponde à execução prospectiva de todas as atividades conforme a abordagem proposta, uma vez que a identificação das realizações e a caracterização das transformações foram sistematizadas posteriormente. Assim, os resultados demonstram a possibilidade de reconstruir e relacionar essas atividades no estudo de caso, mas não que tenham orientado a migração original.

\paragraph{C2 -- Utilidade dos artefatos e resultados.}

A principal contribuição observada está nas relações estabelecidas entre os resultados das diferentes etapas. O conhecimento recuperado na Compreensão de Programas forneceu informações sobre a organização e as interfaces da BSN, utilizadas na interpretação dos comportamentos e na identificação de suas realizações. Os cenários BDD, por sua vez, estabeleceram contratos comportamentais que serviram tanto à construção da referência no ROS~1 quanto à delimitação das responsabilidades examinadas na implementação.

O cenário BDD03 exemplifica essa relação. A especificação da recepção de leituras permitiu orientar a identificação das entidades envolvidas na comunicação entre sensores e \texttt{CentralHub}. O mapeamento resultante relacionou esse contrato às responsabilidades de transmissão e recepção, posteriormente examinadas na transformação T3. Na análise comportamental, o mesmo contrato forneceu uma referência para comparar os oráculos e os resultados de execução nas duas implementações.


De modo semelhante, a investigação das interfaces associadas a BDD15 permitiu relacionar o contrato à responsabilidade de coleta e encaminhamento de registros exercida por \texttt{Probe}. Esse mapeamento forneceu uma referência para caracterizar T2, na qual \texttt{SystemMonitor} mantém essa responsabilidade e incorpora funcionalidades adicionais de supervisão no ROS~2. As correspondências entre responsabilidades e transformações foram consolidadas no registro apresentado na Tabela~\ref{tab:res-migration-record}.

Esses resultados mostram que os artefatos recuperados e os mapeamentos reconstruídos podem ser utilizados conjuntamente para relacionar contratos comportamentais, responsabilidades de implementação e transformações de migração. No estudo retrospectivo, essa utilidade foi observada na organização e interpretação das evidências.

\paragraph{C3 -- Limitações identificadas.}

A aplicação da abordagem depende da disponibilidade de informações sobre o sistema legado. Na BSN, os trabalhos anteriores forneceram documentação e especificações comportamentais que facilitaram a reconstrução. Entretanto, não foram estabelecidos critérios objetivos para determinar a suficiência do conhecimento recuperado ou a abrangência dos contratos selecionados.

A identificação das realizações dependeu da interpretação do código-fonte, dos grafos de chamadas e das interfaces de comunicação. As associações obtidas permitem localizar entidades relacionadas aos comportamentos, mas não constituem uma verificação dinâmica da participação de cada método na execução dos cenários BDD.

Na análise comportamental, os oráculos baseados nas interfaces ROS permitiram examinar comportamentos observáveis nas duas implementações, mas nem sempre estabelecer uma associação inequívoca entre os estímulos e os resultados observados. Essas limitações restringem as conclusões de preservação comportamental, discutidas na RQ2.

Por fim, o caráter retrospectivo impediu a recuperação completa das decisões que orientaram a migração original. Além disso, a ausência de indicadores para avaliar sistematicamente a aplicação do processo limita sua mensurabilidade. 

\paragraph{Síntese da RQ1.}

Em suma, o estudo de caso demonstrou a possibilidade de sistematizar retrospectivamente a migração da BSN por meio das atividades propostas. A abordagem permitiu relacionar especificações comportamentais, evidências da implementação legada, responsabilidades selecionadas, transformações identificadas no ROS~2 e resultados de validação. Essa integração constitui sua principal contribuição observada para a análise do processo de migração.

Entretanto, os resultados sustentam a aplicabilidade da abordagem como procedimento de reconstrução e análise no caso investigado, sem demonstrar sua eficácia na condução prospectiva de uma migração. Sua aplicação em outros sistemas dependerá das evidências disponíveis, da observabilidade dos comportamentos selecionados e da possibilidade de identificar e relacionar suas realizações nas implementações.





\section{RQ2 -- Preservação Comportamental}
\label{sec:avaliacao-rq2}

A RQ2 examina quais evidências recuperadas permitem avaliar a preservação dos contratos comportamentais selecionados durante a migração da BSN. A análise considera a correspondência entre os contratos, os oráculos implementados e os resultados de execução apresentados na Seção~\ref{sec:resultados-validacao-bsn}.

% TODO: Atualizar a avaliação com os resultados das execuções próprias de check_bsn, health_status e dos três casos de BDD15. Enquanto essas execuções não estiverem documentadas, as aprovações aqui discutidas correspondem aos registros históricos selecionados no Capítulo 5.

\paragraph{Correspondência entre os contratos.}

A correspondência entre cenários BDD não implica a equivalência de seus oráculos. Em BDD03, a comparação de valores no ROS~1 incorpora uma restrição temporal que não está presente na asserção examinada do ROS~2. Em BDD05, o teste do legado verifica a ausência de dados na saída capturada do CentralHub, enquanto o teste do ROS~2 verifica a ausência de valores nos campos de risco. Essas diferenças delimitam os aspectos comparáveis dos comportamentos e impedem interpretar a aprovação dos testes como demonstração de equivalência integral dos contratos. A revisão dos oráculos de BDD08 e BDD09 estabelece critérios correspondentes, mas os relatórios históricos consultados não permitem confirmar a execução com essas versões revisadas dos testes.

% TODO: Após as novas execuções de health_status, confirmar as versões dos oráculos e atualizar esta ressalva. Não atribuir os resultados históricos aos oráculos harmonizados sem os respectivos relatórios.

\paragraph{Alcance das evidências de preservação.}

A Tabela~\ref{tab:avaliacao-rq2-contratos} sintetiza as evidências disponíveis e as conclusões sustentadas para os contratos selecionados.

\begin{table}[htbp]
\centering
\caption{Alcance das evidências de preservação comportamental na amostra selecionada}
\label{tab:avaliacao-rq2-contratos}
\small
\begin{tabularx}{\textwidth}{@{}p{0.20\textwidth}XX@{}}
\toprule
\textbf{Contrato} & \textbf{Evidência disponível} & \textbf{Conclusão sustentada} \\
\midrule
BDD03 / BSN-C01 &
Aprovação dos testes de recepção de leituras nas duas implementações, incluindo a revalidação no ROS~2. &
Evidências de preservação da transmissão e recepção de valores nas condições examinadas, sem associação inequívoca entre cada mensagem publicada e recebida. \\
\addlinespace
BDD04 / BSN-P09 &
Aprovação das verificações de correspondência dos valores de risco e presença de risco alto nas interfaces. &
Evidências de preservação da propagação de risco alto, sem demonstração de equivalência dos mecanismos internos de classificação. \\
\addlinespace
BDD05 / BSN-C02 &
Ausência de dados na saída capturada do CentralHub no ROS~1 e ausência de valores nos campos de risco verificados no ROS~2 durante sua indisponibilidade. &
Correspondência quanto à ausência dos resultados de risco nas condições investigadas, sem implicar equivalência de todas as saídas ou dos mecanismos de controle da execução. \\
\addlinespace
BDD08 / BSN-P10 &
Aprovação do cenário de disponibilização do estado de saúde nos registros históricos das duas implementações. &
Evidências de disponibilização do estado de saúde, sem confirmação, nos registros consultados, da execução com os oráculos harmonizados. \\
\addlinespace
BDD09 / BSN-C04 &
Aprovação do cenário de ausência de estado de saúde durante a indisponibilidade do CentralHub nas duas implementações. &
Evidências de ausência dessa informação nas condições examinadas, sujeitas à mesma ressalva sobre os oráculos revisados. \\
\addlinespace
BDD15 / BSN-C07 &
Aprovação das verificações das interfaces de comunicação entre o coletor de registros e o sistema de registro nas duas implementações. &
Evidências de manutenção das relações de comunicação esperadas, sem verificação direta da entrega dos registros ou da supervisão por \textit{heartbeats}. \\
\bottomrule
\end{tabularx}
\end{table}

\begin{comment}
Tabela anterior, preservada para consulta. Inclui cenários e verificações complementares que não integram a amostra atual.
\begin{table}[htbp]
\centering
\caption{Alcance das evidências de preservação comportamental na amostra selecionada}
\label{tab:avaliacao-rq2-contratos}
\small
\begin{tabularx}{\textwidth}{@{}p{0.20\textwidth}XX@{}}
\toprule
\textbf{Contrato} &
\textbf{Evidência disponível} &
\textbf{Conclusão sustentada} \\
\midrule

BDD03 / BSN-C01 &
Aprovação dos testes de recepção de leituras nas duas versões, incluindo revalidação no ROS~2. Os oráculos comparam valores capturados nos tópicos dos sensores e do \texttt{CentralHub}. &
Há evidência de preservação da recepção de valores nas condições examinadas, sem associação inequívoca entre cada leitura publicada e recebida. \\

\addlinespace

BDD04 / BSN-P09 &
Aprovação da observação de risco alto nas duas versões. Quatro testes complementares verificaram a classificação de valores selecionados, com detecção do defeito B7 reintroduzido em uma cópia da implementação C++. &
Há evidência de correspondência na observação de risco alto e na classificação dos valores examinados, sem demonstração de equivalência integral do processamento de risco. \\

\addlinespace

BDD06 / BSN-P08;
BDD07 / BSN-C03 &
Aprovação nominal após revalidação e falhas esperadas nos cenários negativos, com diferenças nas etapas em que ocorreram. Testes controlados recuperaram 30 registros de \texttt{Status} em cada implementação. &
O encaminhamento e a persistência dos registros de \texttt{Status} foram observados, mas a persistência das leituras fisiológicas permanece não demonstrada. \\

\addlinespace

BDD08 / BSN-P10;
BDD09 / BSN-C04 &
Aprovação dos cenários positivo e negativo nas duas implementações, após o alinhamento dos critérios de asserção. &
Há evidência de atendimento aos critérios selecionados de disponibilização do estado de saúde, sem demonstração de equivalência das sequências produzidas. \\

\addlinespace

BDD20 / BSN-C11 &
Falhas esperadas nos cenários BDD. Quatro testes controlados foram aprovados nas duas versões após uma correção no controlador ROS~1. &
Há correspondência nas decisões dos controladores para os valores examinados, considerando a versão corrigida do ROS~1. A recuperação integrada e a frequência efetivamente aplicada não foram demonstradas. \\

\bottomrule
\end{tabularx}
\end{table}
\end{comment}

\begin{comment}
Resultados complementares e cobertura estrutural do bsn_traces, preservados para eventual reinclusão com atribuição e delimitação do escopo.
\paragraph{Evidências complementares e cobertura estrutural.}

Os testes complementares ampliaram a investigação de responsabilidades específicas que não eram suficientemente verificadas pelos cenários BDD. Entre os resultados, destacam-se as verificações controladas da classificação de risco, do processamento do \texttt{CentralHub}, da persistência de registros e das decisões dos controladores de adaptação. Entretanto, suas conclusões permanecem restritas às condições e aos valores examinados.

Além dessas verificações, os relatórios de cobertura permitem examinar a abrangência estrutural das execuções de teste disponíveis. A Tabela~\ref{tab:rq2-cobertura} apresenta os resultados gerais do ROS~1 e os recortes de testes de pacotes, BDD e sua união no ROS~2.

\begin{table}[htbp]
\centering
\caption{Cobertura estrutural das execuções de testes da BSN}
\label{tab:rq2-cobertura}
\small
\begin{tabularx}{\textwidth}{@{}Xrrrr@{}}
\toprule
\textbf{Execução} &
\textbf{Linhas} &
\textbf{Cob. linhas} &
\textbf{Ramos} &
\textbf{Cob. ramos} \\
\midrule

ROS~1 -- Geral (26/10) &
3.443/4.263 & 80,76\% & -- & -- \\

ROS~2 -- Pacotes (01/10) &
2.549/3.147 & 81,00\% & 596/888 & 67,12\% \\

ROS~2 -- BDD (01/10) &
2.368/3.147 & 75,25\% & 530/888 & 59,68\% \\

ROS~2 -- União (01/10) &
2.724/3.147 & 86,56\% & 674/888 & 75,90\% \\

\bottomrule
\end{tabularx}

\par\smallskip
\begin{minipage}{\textwidth}
\footnotesize
Fonte: relatórios de cobertura preservados no repositório
\texttt{bsn\_traces}
(\url{https://github.com/richardjlv/bsn_traces/tree/main/cobertura_geral}).
Os resultados correspondem às execuções e aos recortes de código de produção documentados no repositório.
\end{minipage}
\end{table}

No ROS~2, a união dos testes de pacotes e BDD apresentou cobertura superior à obtida por cada conjunto isoladamente, indicando que as duas suítes exercitaram regiões parcialmente distintas da implementação. Esses resultados consideram conjuntos de testes mais abrangentes do que a amostra de cenários BDD selecionada neste estudo e evidenciam a contribuição da suíte BDD para a cobertura estrutural da implementação migrada.
\end{comment}


\paragraph{Limitações da comparação.}

As conclusões dependem da correspondência dos critérios de aceitação, dos oráculos e das condições de execução. Em particular, a comparação de valores nas interfaces permite verificar a transmissão de leituras, mas não a associação individual entre mensagens; a verificação das relações de comunicação em BDD15 não permite concluir que os registros tenham sido efetivamente encaminhados.

A falha do caso do oxímetro na execução histórica do ROS~2 e sua aprovação posterior exigem distinguir os resultados das diferentes execuções, sem atribuir a falha a uma regressão sem investigação de sua causa. Também é necessário distinguir os resultados históricos de \texttt{health\_status} das execuções com os oráculos revisados. Essas diferenças limitam as conclusões sobre a preservação dos comportamentos selecionados.

\paragraph{Síntese da RQ2.}

Em resposta à RQ2, as evidências recuperadas sustentam conclusões delimitadas sobre a preservação da transmissão e recepção de leituras, da propagação de risco alto e da ausência de resultados de risco durante a indisponibilidade do CentralHub. Para a disponibilização do estado de saúde, os resultados históricos indicam correspondências, cuja interpretação depende da confirmação das versões dos oráculos. Em BDD15, as evidências sustentam a manutenção das interfaces de comunicação associadas ao encaminhamento de registros, mas não a preservação da entrega de seu conteúdo. A aprovação dos casos selecionados, portanto, não demonstra a equivalência comportamental integral das implementações.

\begin{comment}
Síntese anterior, dependente dos cenários e testes complementares retirados do recorte atual.
Em resposta à RQ2, as evidências recuperadas permitiram identificar correspondências comportamentais entre ROS~1 e ROS~2 para os critérios efetivamente verificados. Os resultados sustentam evidências de preservação na recepção de leituras, na observação e classificação de risco para valores selecionados e na disponibilização do estado de saúde. Os testes complementares também forneceram evidências sobre o processamento do \texttt{CentralHub}, a persistência de registros de estado e as decisões dos controladores sob entradas controladas.

Entretanto, a persistência das leituras fisiológicas, a recuperação integrada da adaptação e a equivalência integral dos comportamentos permanecem não demonstradas. Assim, a avaliação permitiu estabelecer conclusões delimitadas sobre a preservação dos contratos selecionados, considerando o alcance dos oráculos e das evidências disponíveis.
\end{comment}

\section{RQ3 -- Caracterização das Transformações de Migração}
\label{sec:avaliacao-rq3}

A RQ3 examina como as responsabilidades de implementação associadas aos contratos comportamentais selecionados foram transformadas durante a migração da BSN do ROS~1 para o ROS~2. A análise parte das responsabilidades delimitadas na Tabela~\ref{tab:res-escopo-identificado} e das transformações T1--T3 apresentadas na Seção~\ref{sec:resultados-planejamento-transformacoes}, examinando a distribuição das responsabilidades, os mecanismos de organização e as características estruturais das classes envolvidas.

\paragraph{Delimitação das responsabilidades analisadas.}

A análise de T1 considera o controle da execução e da atividade dos nós, envolvendo a classe-base \texttt{Component}, \texttt{Sensor} e \texttt{CentralHub} no ROS~1 e a infraestrutura \texttt{RosComponent} e \texttt{LifecycleManager} no ROS~2. Em T2, a responsabilidade de coleta e encaminhamento de informações exercida por \texttt{Probe}, instanciada pelo executável \texttt{collector}, foi relacionada à implementação de \texttt{SystemMonitor} no ROS~2.\footnote{Ponto de entrada de \texttt{collector} no ROS~1: \url{https://github.com/gabrielbfranca/bsn/blob/metrics/src/sa-bsn/target_system/probes/collector/apps/collector.cpp}. Implementação de \texttt{Probe}: \url{https://github.com/gabrielbfranca/bsn/blob/metrics/src/archlib/src/target_system/Probe.cpp}. Implementação de \texttt{SystemMonitor}: \url{https://github.com/rdinizcal/bsn2/blob/metrics/system_monitor/system_monitor/node_monitor.py}.}

Em T3, foram examinadas as responsabilidades de coleta, processamento, comunicação e avaliação de risco dos sensores e do \texttt{CentralHub}, incluindo a disponibilização do estado de saúde do paciente.

\begin{comment}
Recorte anterior com ParamAdapter, associado aos cenários complementares fora da amostra atual.
Em T3, foram examinadas as responsabilidades de coleta, processamento, comunicação e avaliação de risco dos sensores e do \texttt{CentralHub}. A classe \texttt{ParamAdapter} foi incluída como recorte complementar da reconfiguração, relacionada aos contratos BDD17, BDD18 e BDD23, sem constituir uma quarta transformação selecionada.
\end{comment}

As associações aos contratos BDD, registradas na Tabela~\ref{tab:res-migration-record}, fornecem o contexto comportamental para examinar essas responsabilidades. Entretanto, essas relações não demonstram que os contratos tenham orientado as decisões originais de migração nem que as transformações identificadas assegurem a preservação dos comportamentos.

\paragraph{Métricas estruturais das classes selecionadas.}

As Tabelas~\ref{tab:rq3-metricas-target-system-ros1} e~\ref{tab:rq3-metricas-target-system-ros2} apresentam as métricas estruturais das classes examinadas. A complexidade ciclomática (CC) dos métodos foi medida com Lizard~1.26.0 nas duas implementações.\footnote{Saídas nativas do Lizard: ROS~1, \url{https://github.com/gabrielbfranca/bsn/blob/metrics/reports/ros1_static/native/lizard-native.csv}; ROS~2, \url{https://github.com/rdinizcal/bsn2/blob/metrics/reports/class_modularity/lizard-native.csv}.}

% TODO: Conferir se as classes e os relatórios dos branches metrics correspondem às versões de código investigadas nas Seções 5.3 e 5.4, em particular data-persistance-testing no ROS 2. Registrar os commits e, se necessário, atualizar as medições antes de interpretar os valores.

O WMC corresponde à soma das complexidades ciclomáticas dos métodos definidos diretamente em cada classe, enquanto a ``CC máx.'' representa a maior complexidade ciclomática entre esses métodos. Os métodos herdados não são incluídos no WMC da classe derivada. Os valores foram consolidados por scripts a partir das medições, conforme os procedimentos documentados nos repositórios.\footnote{Tabelas consolidadas: ROS~1, \url{https://github.com/gabrielbfranca/bsn/blob/metrics/reports/ros1_static/classes.csv}; ROS~2, \url{https://github.com/rdinizcal/bsn2/blob/metrics/reports/class_modularity/classes.csv}. Procedimentos: \url{https://github.com/gabrielbfranca/bsn/blob/metrics/tools/ros1_metrics/README.md} e \url{https://github.com/rdinizcal/bsn2/blob/metrics/tools/CLASS_MODULARITY.md}.}

\begin{table}[htbp]
\centering
\caption{Complexidade das classes selecionadas na implementação ROS~1}
\label{tab:rq3-metricas-target-system-ros1}
\small
\begin{tabularx}{\textwidth}{@{}Xrr@{}}
\toprule
\textbf{Classe} & \textbf{WMC} & \textbf{CC máx.} \\
\midrule
\multicolumn{3}{@{}l}{\textbf{Infraestrutura de execução (T1)}} \\
\texttt{arch::target\_system::Component} & 22 & 5 \\
\addlinespace
\multicolumn{3}{@{}l}{\textbf{Sensor: coleta e avaliação de risco (T3)}} \\
\texttt{Sensor} & 31 & 7 \\
\texttt{G3T1\_3} & 17 & 4 \\
\texttt{SensorConfiguration} & 33 & 6 \\
\addlinespace
\multicolumn{3}{@{}l}{\textbf{CentralHub: recepção e processamento (T3)}} \\
\texttt{CentralHub} & 21 & 7 \\
\texttt{G4T1} & 51 & 14 \\
\addlinespace
\multicolumn{3}{@{}l}{\textbf{Monitoramento (T2)}} \\
\texttt{arch::target\_system::Probe} & 8 & 1 \\
% \addlinespace
% \multicolumn{3}{@{}l}{\textbf{Reconfiguração (recorte complementar)}} \\
% \texttt{ParamAdapter} & 10 & 3 \\
\bottomrule
\end{tabularx}
\end{table}

\begin{table}[htbp]
\centering
\caption{Complexidade das classes selecionadas na implementação ROS~2}
\label{tab:rq3-metricas-target-system-ros2}
\small
\begin{tabularx}{\textwidth}{@{}Xrr@{}}
\toprule
\textbf{Classe} & \textbf{WMC} & \textbf{CC máx.} \\
\midrule
\multicolumn{3}{@{}l}{\textbf{Infraestrutura compartilhada e ciclo de vida (T1)}} \\
\texttt{RosComponent} & 8 & 2 \\
\texttt{LifecycleManager} & 57 & 7 \\
\texttt{BatteryManager} (compartilhado) & 17 & 8 \\
\addlinespace
\multicolumn{3}{@{}l}{\textbf{Sensor: coleta e avaliação de risco (T3)}} \\
\texttt{Sensor} & 29 & 7 \\
\texttt{DataProcessor} & 18 & 6 \\
\texttt{RiskManager} & 13 & 4 \\
\texttt{RiskEvaluator} & 25 & 7 \\
\texttt{ConfigManager} (sensor) & 3 & 1 \\
\texttt{PublisherManager} (sensor) & 5 & 2 \\
\addlinespace
\multicolumn{3}{@{}l}{\textbf{CentralHub: recepção e processamento (T3)}} \\
\texttt{CentralHub} & 31 & 7 \\
\texttt{SensorDataHandler} & 10 & 7 \\
\texttt{DataFusionEngine} & 22 & 13 \\
\texttt{PublisherManager} (central) & 6 & 3 \\
\addlinespace
\multicolumn{3}{@{}l}{\textbf{Monitoramento (T2)}} \\
\texttt{SystemMonitor} & 33 & 8 \\
% \addlinespace
% \multicolumn{3}{@{}l}{\textbf{Reconfiguração (recorte complementar)}} \\
% \texttt{ParamAdapter} & 15 & 5 \\
\bottomrule
\end{tabularx}
\end{table}

\paragraph{Interpretação das transformações.}

Em T1, o controle da execução e da atividade dos componentes, anteriormente realizado por mecanismos internos de \texttt{Component}, \texttt{Sensor} e \texttt{CentralHub}, passou a utilizar uma infraestrutura compartilhada baseada no ciclo de vida do ROS~2. A classe \texttt{RosComponent} herda de \texttt{LifecycleNode}, enquanto \texttt{LifecycleManager} complementa o gerenciamento dos estados operacionais. Essa mudança caracteriza a adaptação da responsabilidade aos mecanismos da plataforma de destino. As métricas de \texttt{RosComponent} (WMC 8; CC máxima 2) e \texttt{LifecycleManager} (57; 7) descrevem a complexidade das classes, mas não permitem concluir que houve redução da complexidade dessa responsabilidade em relação ao ROS~1, pois sua realização foi distribuída entre diferentes entidades. Os cenários associados à indisponibilidade do CentralHub examinam resultados observáveis, sem verificar diretamente as transições de ciclo de vida.

Em T2, a responsabilidade anteriormente associada a \texttt{Probe} foi mantida e ampliada em \texttt{SystemMonitor}. Além do encaminhamento de eventos, estados e informações de energia, a implementação ROS~2 incorpora a supervisão da atividade dos nós por \textit{heartbeats}. Essa ampliação funcional não decorre necessariamente da substituição do \textit{framework}. BDD15 verifica a correspondência das interfaces de encaminhamento, mas não a entrega dos registros nem a supervisão acrescentada. O WMC de \texttt{Probe} é 8, com CC máxima de 1, enquanto \texttt{SystemMonitor} apresenta WMC de 33 e CC máxima de 8. Entretanto, esses valores não representam uma comparação de funcionalidades equivalentes, pois o componente migrado possui responsabilidades adicionais e o WMC de \texttt{Probe} não inclui os métodos herdados de \texttt{ROSComponent}.

Em T3, a transformação caracteriza-se pela redistribuição das responsabilidades de coleta, processamento, comunicação e avaliação de risco. No ROS~1, \texttt{Sensor::body()} coordena essas operações por meio de métodos virtuais implementados pelas subclasses, como \texttt{G3T1\_3}, caracterizando o padrão \textit{Template Method}. No ROS~2, \texttt{Sensor} mantém a coordenação da execução, mas delega operações a \texttt{DataProcessor}, \texttt{RiskManager} e outros componentes especializados. Da mesma forma, \texttt{CentralHub} utiliza \texttt{SensorDataHandler} e \texttt{DataFusionEngine} para recepção e processamento dos dados. Essa reorganização amplia a utilização de composição e delegação na realização das responsabilidades selecionadas.

As métricas complementam essa caracterização. Nas classes homônimas, o WMC de \texttt{Sensor} passa de 31 para 29 e o de \texttt{CentralHub} de 21 para 31, enquanto a CC máxima permanece em 7. Entretanto, essas diferenças não correspondem à variação da complexidade funcional completa, uma vez que parte das operações passou a ser realizada por classes auxiliares. A comparação entre \texttt{G3T1\_3} e \texttt{G4T1} e seus componentes correspondentes no ROS~2 apresenta a mesma limitação, devido à alteração na distribuição das responsabilidades.

Outro mecanismo identificado na organização do ROS~2 é a utilização dos métodos \path{_create_config_manager()} e \path{_create_publisher_manager()}, declarados em \texttt{RosComponent} e especializados por \texttt{Sensor} e \texttt{CentralHub}. Essa estrutura corresponde ao padrão \textit{Factory Method}, permitindo a criação de gerenciadores específicos a partir de uma infraestrutura compartilhada.

\begin{comment}
Interpretação das métricas do recorte complementar de ParamAdapter, mantida para consulta.
O recorte complementar de \texttt{ParamAdapter} apresenta WMC de 10 no ROS~1 e 15 no ROS~2, com CC máxima de 3 e 5, respectivamente. Esses valores registram características estruturais da implementação de reconfiguração, mas não constituem evidência de uma quarta transformação formalmente caracterizada neste estudo.
\end{comment}

\paragraph{Limites da comparação estrutural.}

As medições de WMC e CC descrevem características das classes selecionadas, mas não representam diretamente a complexidade das responsabilidades completas. A alteração na distribuição das operações e a exclusão dos métodos herdados impedem interpretar as diferenças como uma medida de melhoria ou deterioração da modularidade.

\begin{comment}
Medições exploratórias de coesão e acoplamento. Mantidas para consulta, sem integrar a avaliação ativa em razão das diferenças entre os procedimentos de medição.
Também foram examinadas medições de coesão e acoplamento disponíveis nos repositórios. Na amostra de T3, o LCOM4 bruto reportado pelo \texttt{cpplcom} foi 3 para \texttt{G3T1\_3} e 2 para \texttt{G4T1}, enquanto o \texttt{pyscn} reportou LCOM4 igual a 1 para \texttt{DataProcessor}, \texttt{RiskManager}, \texttt{RiskEvaluator} e \texttt{SensorDataHandler}.\footnote{Resultados nativos de LCOM4: ROS~1, \url{https://github.com/gabrielbfranca/bsn/blob/metrics/reports/ros1_static/native/cpplcom-native.txt}; ROS~2, \url{https://github.com/rdinizcal/bsn2/blob/metrics/reports/lcom4/pyscn.json}.} Entretanto, esses valores não são diretamente comparáveis devido às diferenças nos critérios utilizados pelas ferramentas. A mesma restrição se aplica ao CBO, calculado pelo CCCC no ROS~1 e pelo \texttt{pyscn} no ROS~2.\footnote{Resultados nativos de acoplamento: ROS~1, \url{https://github.com/gabrielbfranca/bsn/blob/metrics/reports/ros1_static/native/cccc/cccc.xml}; ROS~2, \url{https://github.com/rdinizcal/bsn2/blob/metrics/reports/class_modularity/pyscn.json}.}
\end{comment}

Dessa forma, a avaliação estrutural permite caracterizar a distribuição das responsabilidades e os mecanismos utilizados nas implementações, mas não demonstrar uma melhoria quantitativa global da qualidade estrutural do sistema.

\paragraph{Síntese da RQ3.}

Em resposta à RQ3, a análise permitiu caracterizar três formas de transformação das responsabilidades associadas aos contratos comportamentais selecionados: adaptação tecnológica do controle de execução ao modelo de ciclo de vida do ROS~2 (T1), manutenção e ampliação funcional do monitoramento (T2) e redistribuição das operações de processamento, comunicação e avaliação de risco entre componentes especializados (T3).

A inspeção das implementações e as métricas estruturais complementaram a caracterização dessas mudanças, enquanto as relações com os contratos BDD permitiram situá-las em relação aos comportamentos investigados. Os resultados descrevem as transformações identificadas retrospectivamente, sem demonstrar que a abordagem proposta tenha orientado as decisões originais da migração ou que as mudanças estruturais impliquem, por si só, preservação comportamental.

\section{Ameaças à Validade}

O caráter retrospectivo do estudo constitui uma limitação para a avaliação da abordagem proposta, uma vez que nem todas as atividades e decisões da migração original puderam ser reconstruídas a partir das evidências disponíveis. Embora tenha sido possível relacionar práticas realizadas, artefatos recuperados e procedimentos sistematizados posteriormente, os resultados não demonstram que a metodologia tenha orientado originalmente a migração, nem permitem avaliar seu impacto sobre a migração.

A avaliação também foi limitada a uma amostra de comportamentos e responsabilidades da BSN. Na análise comportamental, as conclusões dependem dos critérios efetivamente verificados pelos oráculos e da comparabilidade das condições de execução entre ROS~1 e ROS~2. A predominância de testes orientados às interfaces de comunicação permite observar resultados sem depender da organização interna dos componentes, mas nem sempre estabelece uma correspondência inequívoca entre estímulos e respostas. Dessa forma, as evidências obtidas sustentam conclusões sobre os contratos selecionados, sem demonstrar a preservação comportamental integral do sistema. A identificação das realizações, baseada na inspeção do código-fonte, nos grafos de chamadas e nas interfaces, não permite confirmar dinamicamente a participação de cada entidade na execução dos cenários.

Na caracterização das transformações estruturais, as diferenças de linguagem, decomposição das responsabilidades e critérios de medição limitam a comparação quantitativa entre as implementações. Assim, as métricas utilizadas complementam a interpretação das transformações identificadas, mas não demonstram isoladamente uma melhoria global da modularidade.

